\section{หลักการและเหตุผล}

จังหวัดขอนแก่นนับเป็นจังหวัดที่มีศักยภาพทางการท่องเที่ยวสูง ด้วยความหลากหลายของแหล่งท่องเที่ยว ทั้งทางธรรมชาติ วัฒนธรรม ประวัติศาสตร์ ตลอดจนแหล่งท่องเที่ยวเชิงวิชาการที่โดดเด่นอย่าง ``ไดโนเสาร์'' ซึ่งถือเป็นอัตลักษณ์ที่สร้างชื่อเสียงให้กับจังหวัด ขอนแก่นจึงเป็นหนึ่งในจุดหมายสำคัญของนักท่องเที่ยวทั้งชาวไทยและต่างชาติที่ต้องการสัมผัสประสบการณ์ที่หลากหลายและแตกต่าง

อย่างไรก็ตาม ปัญหาสำคัญที่ยังคงปรากฏอยู่คือ การเข้าถึงข้อมูลการท่องเที่ยวที่กระจัดกระจายอยู่ในหลายช่องทาง ไม่ว่าจะเป็นเว็บไซต์ หน่วยงานราชการ เพจโซเชียลมีเดีย หรือแอปพลิเคชันทั่วไป นักท่องเที่ยวจำเป็นต้องใช้เวลาค้นหาข้อมูลจากหลายแหล่ง ซึ่งมักพบปัญหาว่าข้อมูลไม่ครบถ้วน ไม่ทันสมัย หรือขาดการรองรับภาษาสำหรับนักท่องเที่ยวต่างชาติ ส่งผลให้การวางแผนท่องเที่ยวไม่มีประสิทธิภาพและอาจทำให้นักท่องเที่ยวพลาดโอกาสในการเข้าถึงประสบการณ์ที่ควรได้รับ

เพื่อแก้ไขข้อจำกัดดังกล่าว จึงเกิดแนวคิดในการพัฒนาโครงงาน ``ไดโน : เว็บแอปพลิเคชันศูนย์กลางของการท่องเที่ยวจังหวัดขอนแก่น'' ซึ่งเป็นแอปพลิเคชันที่มุ่งเน้นการรวมศูนย์ข้อมูลการท่องเที่ยวอย่างครบวงจร ภายในแอปประกอบด้วยฟังก์ชันการค้นหาแหล่งท่องเที่ยว ร้านอาหาร ที่พัก และกิจกรรมต่าง ๆ พร้อมข้อมูลรีวิวจากผู้ใช้งานจริง รวมถึงระบบแนะนำเส้นทางและระบบแชตบอต (Chatbot) ที่สามารถโต้ตอบให้คำแนะนำได้แบบเรียลไทม์ ช่วยให้การค้นหาข้อมูลและการวางแผนทริปเป็นไปอย่างสะดวก รวดเร็ว และทันสมัย

นอกจากการอำนวยความสะดวกแก่นักท่องเที่ยวแล้ว แอปพลิเคชันยังเป็นเครื่องมือสำคัญในการสนับสนุนผู้ประกอบการท้องถิ่นให้สามารถโปรโมทธุรกิจได้ตรงกลุ่มเป้าหมาย อีกทั้งยังสามารถเก็บสถิติการใช้งาน เช่น จุดท่องเที่ยวยอดนิยม ช่วงฤดูกาลท่องเที่ยว หรือกิจกรรมที่ได้รับความสนใจสูง ซึ่งข้อมูลเหล่านี้สามารถนำไปใช้วิเคราะห์แนวโน้มและวางแผนการพัฒนาการท่องเที่ยวของจังหวัดในอนาคต

ดังนั้น การพัฒนาแอปพลิเคชัน ``ไดโน'' จึงไม่เพียงเป็นการแก้ปัญหาด้านการเข้าถึงข้อมูลการท่องเที่ยวเท่านั้น แต่ยังมีบทบาทสำคัญในการสร้างประสบการณ์การท่องเที่ยวที่ครบถ้วนสำหรับผู้มาเยือน และสนับสนุนเศรษฐกิจท้องถิ่นอย่างยั่งยืน ทั้งยังเป็นการวางรากฐานที่มั่นคงต่อการพัฒนาการท่องเที่ยวของจังหวัดขอนแก่นในระยะยาว
